\chapter{Breast Uplift (Mastopexy)}

\section*{What Is This Surgery?}
Mastopexy, commonly known as a breast lift, is a surgical procedure designed to elevate and reshape sagging (ptotic) breasts. This involves the removal of excess skin, tightening of the underlying breast tissue, and repositioning of the nipple-areola complex (NAC) to a higher, more aesthetically pleasing and youthful position on the breast mound. The procedure addresses breast ptosis, a condition characterized by the descent of the breast tissue and NAC below the inframammary fold, often resulting from factors such as aging, significant weight fluctuations, pregnancy, breastfeeding, and gravitational effects. Beyond aesthetic improvement, mastopexy can alleviate physical discomforts associated with severe ptosis, such as skin irritation in the inframammary fold, and improve the fit of clothing and brassieres, thereby enhancing a patient's quality of life and self-confidence.

\section*{Alternative Names}
\begin{itemize}
  \item Breast Lift Surgery
  \item Breast Ptosis Correction
  \item Nipple Repositioning Surgery
  \item Breast Reshaping Surgery
  \item Mammary Gland Suspension
\end{itemize}

\section*{Relevant Surgical Anatomy}
The breast is a modified apocrine gland situated on the anterior chest wall, overlying the pectoralis major muscle, extending from the second to the sixth rib vertically and from the sternal edge to the mid-axillary line horizontally. It is composed of glandular tissue (15-20 lobes), adipose tissue, and fibrous connective tissue, including Cooper's ligaments, which provide structural support by connecting the breast parenchyma to the overlying skin and deep fascia. The skin envelope, particularly its elasticity and integrity, plays a crucial role in breast shape and position. The nipple-areola complex (NAC) is centrally located, typically at the apex of the breast mound, and its position relative to the inframammary fold (IMF) is a key determinant of breast ptosis. The IMF is a critical surgical landmark, representing the natural crease where the breast meets the chest wall.

The arterial supply to the breast is rich and originates primarily from the internal mammary artery (medial perforators), the lateral thoracic artery (lateral aspect), and the thoracoacromial artery (superior aspect), with contributions from intercostal perforators. Venous drainage largely parallels the arterial supply. The innervation of the breast, particularly for sensation to the NAC, is primarily derived from the lateral cutaneous branches of the intercostal nerves (T4-T6). Preservation of these neurovascular bundles is paramount during mastopexy to maintain NAC viability and sensation. Lymphatic drainage is extensive, predominantly to the axillary lymph nodes (levels I, II, III), but also to internal mammary nodes, which is highly relevant in breast cancer surgery but less directly involved in elective mastopexy unless tissue is sent for pathology.

\section{Surgical Principle \& Rationale}
The fundamental principle of mastopexy involves the creation of a new, more superiorly located and aesthetically pleasing breast mound by excising redundant skin, reshaping the glandular tissue, and repositioning the nipple-areola complex (NAC) while meticulously preserving its vascularity and innervation. The rationale is to counteract the effects of gravity, aging, and other factors that lead to breast ptosis, restoring a youthful contour and projection. The surgeon strategically removes excess skin from the lower pole and periphery of the breast, effectively lifting the entire breast envelope. Concurrently, the underlying breast parenchyma is often reshaped and plicated (tightened with sutures) to create a firmer, more conical, and projected breast contour, providing internal support for the newly elevated breast.

The NAC is then transposed to its new, elevated position, typically at the apex of the newly formed breast mound, ensuring its viability through a carefully designed pedicle of breast tissue (e.g., superior, superomedial, inferior, or central pedicle). Various surgical patterns are employed, such as the periareolar (donut), vertical (lollipop), and inverted T (anchor) techniques, each chosen based on the degree of ptosis, desired outcome, and acceptable scar burden. The overarching technical goals include elevating the NAC above the inframammary fold, reducing areolar diameter if necessary, reshaping and tightening the glandular tissue, excising excess skin to create a tauter breast envelope, and achieving optimal symmetry between both breasts with the least conspicuous scarring possible.

\section{History \& Surgical Evolution}
The evolution of mastopexy techniques is closely intertwined with the history of breast reduction surgery, with early attempts focusing primarily on skin excision. One of the earliest documented procedures for breast ptosis correction was performed by Joseph Thorek in the 1920s, involving significant skin removal and nipple grafting, which often resulted in loss of nipple sensation and viability due to the complete detachment of the NAC. In 1928, Biesenberger described a technique that involved glandular resection and reshaping, marking a crucial shift towards addressing the underlying breast tissue in addition to the skin envelope, thereby improving the durability of the lift.

Throughout the mid-20th century, surgeons like Marc Dufourmentel and Ivo Pitanguy refined techniques, introducing various skin patterns and glandular suspension methods that aimed to preserve nipple sensation and improve aesthetic outcomes. Pitanguy's technique, developed in the 1960s, utilized an inverted T-scar pattern and glandular reshaping, becoming a widely adopted standard for moderate to severe ptosis. The late 20th century saw further advancements, notably with Madeleine Lejour's introduction of the vertical mastopexy in the 1990s, which significantly reduced the horizontal scar in the inframammary fold, leading to less visible scarring and often improved breast projection. Modern mastopexy continues to evolve, emphasizing minimal scarring, long-term stability, and the integration of techniques with breast augmentation or reduction for comprehensive aesthetic improvement, often incorporating internal support matrices or mesh for enhanced durability.

\section{Clinical Indications}
Mastopexy is primarily indicated for individuals experiencing breast ptosis, which can manifest in various degrees and forms, impacting both aesthetics and physical comfort.

\begin{itemize}
  \item Correction of true breast ptosis, where the nipple-areola complex (NAC) falls below the inframammary fold, often classified using systems like the Regnault classification (Grade I, II, III).
  \item Improvement of pseudoptosis, a condition where the breast tissue sags but the nipple remains above the inframammary fold, typically due to glandular atrophy and skin laxity.
  \item Addressing breast asymmetry caused by differential sagging, volume loss, or developmental variations between the breasts.
  \item Repositioning of a malpositioned, elongated, or enlarged nipple-areola complex to a more central and aesthetically appropriate location.
  \item Enhancing breast shape, projection, and firmness following significant weight loss, which often results in deflated and sagging breasts.
  \item Restoring a more youthful and aesthetically pleasing breast contour after pregnancy and breastfeeding, which can lead to glandular involution and skin laxity.
  \item Alleviating physical discomforts such as skin irritation, rashes, or hygiene issues in the inframammary fold due to skin-on-skin contact.
  \item As an adjunct to breast augmentation (mastopexy-augmentation) for patients desiring both a lift and increased breast volume, or to breast reduction (reduction-mastopexy) for individuals with both macromastia and significant ptosis.
\end{itemize}

\section{Clinical Positioning}
Mastopexy serves primarily as a first-line, elective cosmetic procedure for individuals seeking to correct breast ptosis and improve breast aesthetics. It is considered a primary intervention when the main goal is to lift and reshape sagging breasts without significant changes in breast volume, although the improved projection can create an appearance of increased volume. Patients typically present with concerns about breast shape, nipple position, and overall breast youthfulness, making mastopexy the direct solution for these aesthetic goals.

However, mastopexy can also function as a secondary or adjunct procedure in several contexts. It is frequently combined with breast augmentation (mastopexy-augmentation) for patients who desire both a lift and increased breast volume, addressing both ptosis and hypomastia simultaneously. Conversely, it is often performed in conjunction with breast reduction (reduction-mastopexy) for individuals with both macromastia and significant ptosis, where excess breast tissue is removed, and the remaining breast is lifted and reshaped. Furthermore, mastopexy may be considered a secondary or revisionary procedure for patients who have experienced recurrence of ptosis after a previous breast lift, or to correct unsatisfactory outcomes from prior breast surgeries, such as persistent asymmetry, nipple malposition, or undesirable scar formation. In these scenarios, it acts as a corrective or refining intervention.

\section{Surgical Technique \& Procedure}
Mastopexy is typically performed under general anesthesia, ensuring the patient's comfort and immobility throughout the procedure. The patient is positioned supine, with arms extended on arm boards, allowing for full access to the breast and axillary regions. Preoperative markings are crucial and are performed with the patient in an upright position to accurately assess the degree of ptosis, plan the new nipple-areola complex (NAC) position, and delineate the precise skin excision pattern.

The surgical steps generally follow a structured sequence:
\begin{enumerate}
  \item \textbf{Anesthesia and Markings}: General anesthesia is administered. Precise preoperative markings are made on the breast to guide skin excision and NAC repositioning, often using a specialized ruler or template to ensure symmetry and optimal placement.
  \item \textbf{Infiltration}: A local anesthetic solution containing epinephrine (e.g., lidocaine with epinephrine) is infiltrated into the planned incision lines and breast tissue. This serves to minimize intraoperative bleeding through vasoconstriction and provides immediate postoperative pain relief.
  \item \textbf{Skin Excision}: Incisions are made according to the chosen mastopexy pattern (e.g., periareolar, vertical, or inverted T). The predetermined excess skin is carefully excised, creating the new breast envelope and defining the boundaries for glandular reshaping.
  \item \textbf{NAC Transposition}: The NAC is carefully dissected, typically on a pedicle of underlying breast tissue (e.g., superior, superomedial, inferior, or central pedicle), and transposed to its new, higher position. The integrity of the pedicle, which supplies blood and sensation to the NAC, is meticulously preserved.
  \item \textbf{Glandular Reshaping and Plication}: The remaining breast glandular tissue is reshaped and often plicated (folded and sutured) to create a more conical, projected, and firm breast mound. This internal support is critical for maintaining the long-term lift and shape.
  \item \textbf{Hemostasis}: Meticulous hemostasis is achieved throughout the procedure using electrocautery to minimize postoperative bleeding and the risk of hematoma formation.
  \item \textbf{Closure}: The skin flaps are redraped and closed in layers. Deep dermal sutures (often absorbable) are used to approximate the skin edges and provide tension relief, followed by skin closure with fine absorbable or non-absorbable sutures. Drains, typically small suction drains, may be placed temporarily to collect any fluid accumulation (seroma or hematoma).
  \item \textbf{Dressings}: Sterile dressings are applied to the incision sites, and a surgical bra or compression garment is fitted to provide support, reduce swelling, and help mold the new breast shape.
\end{enumerate}
The duration of the procedure typically ranges from 2 to 4 hours, depending on the complexity of the ptosis and whether it is combined with other procedures like augmentation.

\section*{Preoperative Workup \& Preparation}
Thorough preoperative preparation is essential for a safe and successful mastopexy, minimizing risks and optimizing outcomes. This process begins with a comprehensive consultation with the plastic surgeon, where the patient's medical history, surgical goals, and expectations are discussed in detail, including an assessment of breast ptosis severity and skin quality.

Key preparatory steps include:
\begin{itemize}
  \item \textbf{Medical Evaluation}: A complete medical history and physical examination are performed to assess overall health, identify any chronic conditions (e.g., diabetes, hypertension, thyroid disorders), and review all current medications, including over-the-counter drugs, herbal supplements, and vitamins.
  \item \textbf{Laboratory Tests}: Routine preoperative blood tests (e.g., complete blood count, coagulation profile, metabolic panel, pregnancy test) are typically ordered to ensure the patient is fit for surgery.
  \item \textbf{Imaging}: A baseline mammogram or breast ultrasound may be required, especially for patients over 40 years of age or with a family history of breast cancer, to rule out any underlying breast pathology.
  \item \textbf{Medication Adjustments}: Patients are usually advised to stop taking blood-thinning medications (e.g., aspirin, NSAIDs, warfarin, certain herbal supplements like Ginkgo Biloba or Vitamin E) at least two weeks prior to surgery to minimize bleeding risk.
  \item \textbf{Smoking Cessation}: Smoking significantly impairs wound healing, increases the risk of skin necrosis, infection, and other complications. Patients are strongly advised to cease smoking for at least 4-6 weeks before and after surgery.
  \item \textbf{NPO Status}: Patients must adhere to "nil per os" (NPO) guidelines, meaning no food or drink for a specified period (typically 6-8 hours) before surgery, to prevent aspiration during anesthesia.
  \item \textbf{Prophylactic Antibiotics}: A single dose of intravenous prophylactic antibiotics is typically administered preoperatively to reduce the risk of surgical site infection.
  \item \textbf{Informed Consent}: A detailed informed consent form is reviewed and signed, outlining the procedure, potential risks, benefits, alternatives, and expected outcomes, including scar patterns and potential changes in sensation.
  \item \textbf{Preoperative Photography}: Standardized photographs are taken for medical record purposes and to document the preoperative breast appearance for comparison.
\end{itemize}
Patients should arrange for transportation home after surgery and have assistance for the first 24-48 hours of recovery.

\section*{Contraindications \& Precautions}
Careful patient selection is paramount for mastopexy to ensure safety and optimize outcomes.

\textbf{Absolute Contraindications:}
\begin{itemize}
  \item Active breast infection or inflammation (e.g., mastitis, cellulitis).
  \item Untreated or active breast malignancy; any suspicious breast lesions must be fully investigated and treated prior to elective mastopexy.
  \item Uncontrolled systemic medical conditions (e.g., severe heart disease, uncontrolled diabetes, severe autoimmune disorders, severe pulmonary disease) that significantly increase surgical and anesthetic risks.
  \item Severe coagulopathy or bleeding disorders that cannot be adequately managed or corrected.
  \item Pregnancy or active breastfeeding, as breast tissue is hormonally active and results would be unpredictable.
  \item Body dysmorphic disorder or unrealistic expectations regarding surgical outcomes, which may lead to patient dissatisfaction despite a technically successful procedure.
  \item Inability to provide informed consent due to cognitive impairment or legal status.
\end{itemize}

\textbf{Relative Precautions:}
\begin{itemize}
  \item \textbf{Smoking}: While not an absolute contraindication, active smoking significantly increases the risk of wound healing complications, skin necrosis, infection, and delayed recovery. Patients are strongly advised to quit for at least 4-6 weeks pre- and post-operatively.
  \item \textbf{Obesity}: Significant obesity (BMI >30-35) can increase surgical risks, including infection, wound dehiscence, seroma formation, and anesthetic complications. Weight optimization is often recommended.
  \item \textbf{Planned Future Pregnancy/Breastfeeding}: While mastopexy does not typically preclude future breastfeeding, subsequent pregnancies can cause recurrence of ptosis due to hormonal changes and breast engorgement, potentially necessitating revision surgery. This should be thoroughly discussed.
  \item \textbf{Certain Medications}: Chronic use of corticosteroids, immunosuppressants, or certain anticoagulants may require careful management, temporary cessation, or alternative strategies to mitigate risks.
  \item \textbf{Unstable Weight}: Patients with significant weight fluctuations should be advised to achieve a stable weight for at least 6-12 months before surgery to optimize long-term results and minimize recurrence of ptosis.
  \item \textbf{Prior Breast Surgery}: Previous breast surgeries (e.g., augmentation, reduction, biopsy) may alter anatomy, vascularity, and scar tissue, requiring a more cautious approach and potentially limiting surgical options or increasing complexity.
  \item \textbf{Large Breast Volume}: For very large breasts with significant ptosis, a reduction-mastopexy might be a more appropriate procedure to achieve a stable and aesthetically pleasing outcome.
\end{itemize}

\section*{Complications \& Risks}
As with any surgical procedure, mastopexy carries inherent risks, which can be broadly categorized into general surgical risks and procedure-specific complications. Patients must be thoroughly informed of these possibilities during the consent process.

\textbf{General Surgical Risks:}
\begin{itemize}
  \item \textbf{Anesthesia Complications}: Reactions to general anesthesia, including nausea, vomiting, allergic reactions, respiratory depression, or, rarely, more severe cardiovascular events such as myocardial infarction or stroke.
  \item \textbf{Bleeding}: Postoperative hematoma (collection of blood under the skin) requiring drainage, which can increase swelling and pain.
  \item \textbf{Infection}: Surgical site infection, which may require antibiotics or, in severe cases, surgical drainage and debridement.
  \item \textbf{Deep Vein Thrombosis (DVT) and Pulmonary Embolism (PE)}: Formation of blood clots in the legs that can travel to the lungs, a rare but serious complication. Prophylactic measures are typically employed.
  \item \textbf{Poor Wound Healing}: Delayed healing, wound dehiscence (opening of the incision), or skin necrosis, particularly in smokers or patients with compromised circulation.
\end{itemize}

\textbf{Procedure-Specific Risks:}
\begin{itemize}
  \item \textbf{Scarring}: All mastopexy procedures result in permanent scars. While efforts are made to minimize their visibility, hypertrophic (raised, red) or keloid (thick, itchy, extending beyond the wound boundaries) scars can occur, especially in predisposed individuals.
  \item \textbf{Asymmetry}: Despite careful planning and meticulous technique, some degree of asymmetry in breast shape, size, or nipple position may persist or develop postoperatively.
  \item \textbf{Nipple-Areola Complex (NAC) Sensation Changes}: Temporary or permanent changes in nipple sensation, including numbness, hypersensitivity, or complete loss of sensation, are common due to nerve manipulation.
  \item \textbf{Nipple Necrosis}: A rare but serious complication where the blood supply to the NAC is compromised, leading to tissue death and potential partial or complete loss of the nipple. This risk is higher with larger transpositions or in smokers.
  \item \textbf{Contour Irregularities}: Dimpling, puckering, or unevenness of the breast contour, which may require revision surgery.
  \item \textbf{Recurrence of Ptosis}: Over time, due to gravity, aging, weight changes, or subsequent pregnancies, some degree of breast sagging may recur, potentially necessitating revision surgery.
  \item \textbf{Fat Necrosis}: Hard lumps in the breast tissue caused by the death of fat cells, which may be palpable and sometimes require observation or excision.
  \item \textbf{Seroma}: Collection of clear fluid under the skin, which may require aspiration with a needle.
  \item \textbf{Persistent Pain}: Chronic pain in the breast area, though rare, can occur.
  \item \textbf{Unsatisfactory Aesthetic Outcome}: Despite technical success, the patient may not be fully satisfied with the final appearance, highlighting the importance of realistic preoperative expectations.
\end{itemize}

\section*{Postoperative Recovery Timeline}
Immediately after mastopexy, patients are transferred to a Post-Anesthesia Care Unit (PACU) for close monitoring as they recover from anesthesia. Pain medication will be administered intravenously or orally to manage discomfort, and anti-nausea medication may be given if needed. Patients will typically have dressings applied to their breasts and will be fitted with a surgical support bra or compression garment, which must be worn continuously for several weeks to minimize swelling, provide support, and help mold the new breast shape. Drains, if placed, will be monitored for output, and patients are usually discharged home the same day or the following morning, depending on their recovery and the extent of the surgery.

During the first 24-48 hours, it is normal to experience moderate pain, swelling, and bruising, which can be effectively managed with prescribed oral pain relievers. Activity will be restricted, with emphasis on rest and avoiding strenuous movements or lifting. Patients are usually advised to sleep on their back with their upper body elevated to reduce swelling. Over the first 1-2 weeks, swelling and bruising will gradually subside. Drains, if present, are typically removed within the first week during a follow-up visit. Patients are usually advised to avoid showering for a few days and then to gently wash the incision sites, avoiding scrubbing. Strenuous activities, heavy lifting, and upper body exercises are restricted for at least 4-6 weeks to allow for proper tissue healing and prevent complications. Long-term recovery involves scar maturation, which can take 6-12 months, during which scars will gradually soften, flatten, and fade. The final breast shape and position will become apparent as swelling resolves over several months.

\section*{Surgical Findings \& Pathology}
During a mastopexy procedure, the surgeon meticulously documents several key findings. These include the precise degree of breast ptosis (often graded), the quality and elasticity of the skin envelope, the volume and consistency of the glandular and adipose tissue, and the original position and size of the nipple-areola complex (NAC). Any unexpected findings, such as palpable masses, cysts, or areas of induration, are noted and may prompt intraoperative biopsy or further investigation. The surgeon also assesses the vascularity of the breast tissue and the planned pedicle for the NAC to ensure its viability.

For a pure mastopexy, where only skin and minimal breast tissue are excised for reshaping, the resected tissue is typically sent for routine pathological examination only if there was a preoperative concern or an unexpected intraoperative finding. If a significant amount of breast tissue is removed, as in a reduction-mastopexy, the pathology report will detail the histological composition of the resected tissue. Common benign findings include fibrocystic changes, adenosis, duct ectasia, or fat necrosis. The pathologist confirms the absence of malignancy and provides a detailed description of any benign breast disease, which is crucial for the patient's long-term breast health management.

\section{Expected Postoperative Course}
A normal and successful postoperative course following mastopexy is characterized by a series of positive changes and a smooth recovery trajectory. Immediately after surgery, patients will experience a significant lift and reshaping of their breasts, though initial swelling and bruising will be present. Pain is typically moderate in the first few days and gradually resolves with oral pain medication. Swelling and bruising will progressively diminish over the first 2-4 weeks. The incision lines will initially appear red and slightly raised, but with proper care, they are expected to flatten and fade over 6 to 12 months, eventually becoming less noticeable.

The nipple-areola complex (NAC) will be positioned appropriately at or slightly above the inframammary fold, and its sensation may be altered, often improving over several months. Patients can expect to gradually resume light activities within 1-2 weeks and return to most normal activities, including light exercise, by 4-6 weeks. The final, stable breast shape and projection will become evident as all residual swelling resolves and the tissues settle, typically around 3-6 months post-surgery. A successful course also involves the absence of significant complications such as infection, hematoma, or nipple necrosis, leading to high patient satisfaction with their improved breast contour and enhanced self-confidence.

\section{Postoperative Care \& Follow-up}
Postoperative care for mastopexy involves a series of follow-up visits and specific instructions to ensure proper healing and monitor for any complications. Adherence to these guidelines is crucial for optimal results.

\begin{itemize}
  \item \textbf{Initial Postoperative Visits}: Typically scheduled within the first week to remove drains (if placed), check incision sites for signs of infection or dehiscence, and assess early healing.
  \item \textbf{Suture Removal/Care}: External sutures, if non-absorbable, are usually removed at 1-2 weeks. Patients are instructed on proper wound care, including gentle cleansing and keeping incisions dry.
  \item \textbf{Activity Resumption}: Gradual return to normal activities, with specific guidance on when to resume light exercise (usually 3-4 weeks) and strenuous activities or heavy lifting (6-8 weeks).
  \item \textbf{Support Garment Use}: Continued use of a surgical bra or compression garment for several weeks, as advised by the surgeon, to support the breasts, reduce swelling, and aid in shaping.
  \item \textbf{Scar Management}: Recommendations for scar massage, silicone sheets, or topical creams to optimize scar appearance, typically starting a few weeks post-op once incisions are fully closed.
  \item \textbf{Long-term Follow-up}: Subsequent visits at 3, 6, and 12 months to assess long-term results, scar maturation, and address any concerns or potential for revision.
  \item \textbf{Routine Breast Screening}: Patients should continue with their age-appropriate routine breast cancer screening (e.g., mammograms) as recommended by their primary care physician.
\end{itemize}
Patients are advised to contact their surgeon immediately if they experience warning signs such as fever, excessive or worsening pain, significant redness or warmth around the incisions, pus-like discharge, sudden changes in breast shape or sensation, or shortness of breath.

\section*{Epidemiology}
Breast ptosis is a common condition affecting a significant portion of the adult female population, with its prevalence increasing with age, parity (number of pregnancies), and body mass index. While precise incidence rates for mastopexy alone are challenging to isolate due to its frequent combination with augmentation or reduction, it is consistently ranked among the most commonly performed aesthetic breast procedures globally. The typical patient undergoing mastopexy is a woman in her 30s to 50s, often multiparous, who has experienced breast sagging due to pregnancy, breastfeeding, significant weight loss, or the natural aging process.

The demand for mastopexy has steadily risen over the past decades as women seek to restore a more youthful breast contour and improve body image. Mortality rates associated with mastopexy are exceedingly low, consistent with other elective cosmetic surgeries performed under general anesthesia, typically reported to be less than 0.01\%. Morbidity, while generally low, can include complications such as hematoma (1-5\%), infection (1-3\%), and wound healing issues (e.g., dehiscence, skin necrosis), with rates varying based on patient factors (e.g., smoking, obesity, diabetes) and surgical technique. Long-term complications like recurrence of ptosis (5-15%) or changes in nipple sensation (up to 20-30% temporary, 5-10% permanent) are also reported, influencing patient satisfaction and potential for revision surgery.

\section*{Outcomes \& Prognosis}
The outcomes of mastopexy are generally highly favorable, with high rates of patient satisfaction reported in the literature, often exceeding 85-90\%. Studies consistently show that the procedure significantly improves breast shape, projection, and position, leading to enhanced body image, self-esteem, and quality of life. Functional outcomes include alleviation of physical discomforts such as skin irritation in the inframammary fold and improved fit of clothing. Survival outcomes are not directly applicable to mastopexy as it is an elective cosmetic procedure, but the long-term safety profile is excellent.

The durability of the results is influenced by several prognostic factors, including the patient's skin quality and elasticity, breast size, age, and future life events such as pregnancy or significant weight fluctuations. While mastopexy provides a long-lasting lift, the natural aging process and gravity will continue to exert their effects over time, meaning some degree of recurrent ptosis is possible years after the procedure, with reported recurrence rates varying from 5% to 15% over 5-10 years. Prognostic factors for optimal long-term outcomes include meticulous surgical technique, appropriate patient selection (e.g., non-smokers, stable weight), adherence to postoperative care instructions, and a healthy lifestyle. While most patients achieve their desired aesthetic goals with a single procedure, a small percentage may opt for revision surgery to address minor asymmetries, improve scar appearance, or correct recurrent sagging. Overall, mastopexy offers a predictable and effective solution for breast ptosis, with a profoundly positive impact on patients' physical comfort and psychological well-being.

\section*{References}
\begin{enumerate}
  \item MedlinePlus. Breast lift. U.S. National Library of Medicine; 2024. Available from: \url{https://medlineplus.gov/ency/article/002980.htm}
  \item Grabb \& Smith's Plastic Surgery. 8th ed. Thorne CH, Chung KC, Gosain AK, et al., editors. Wolters Kluwer; 2019.
  \item American Society of Plastic Surgeons. Breast Lift. Available from: \url{https://www.plasticsurgery.org/cosmetic-procedures/breast-lift}
  \item Mathes and Hentz Plastic Surgery. 3rd ed. Hentz VR, Mathes SJ, editors. Elsevier Saunders; 2017.
\end{enumerate}

\clearpage